\documentclass[12pt]{article}

\usepackage[a4paper, total={6in, 8in}]{geometry}
\usepackage{textcomp}

\begin{document}
\setlength{\parindent}{0pt}

\def \t {\textrightarrow}
\def \v {\vspace{1em}}
\def \bi {\begin{itemize}}
\def \ei {\end{itemize}}
\def \s[#1] {\section*{#1}}
\def \ss[#1] {\subsection*{#1}}
\def \sss[#1] {\subsubsection*{#1}}

\s[Il romanzo]
\ss[Moravia]
Tradizione manzoniana abbandonata \t poi il disfacimento del romanzo \\
Codice di Perla, Palazzeschi \\
Il romanzo viene letto come esperienza di vita \t durante i poteri forti, con l esperienza della guerra è centrale \\
Fino agli 20 con Tozzi \t decisivi sono gli anni 30 \t 1929 c'è lo spartiacque, con il declino della borghesia e il senso di disfacimento valoriale \t con Morava, che pubblica "Gli indifferenenti"

\ss[Gli indifferenenti]
Parla di una famiglia romana \t il genere descrive la realta del 29, in cui si verificano gli eventi piu dequalificanti \\
Una famiglia in disfacimento economico ma non solo \t conservano il titolo, ma c'è una totale negazione della sostanza interiore che aveva descritto i valori borghesi \\
Decadimento dei personaggi, e descrizione che porta a descrivere il dramma interiore, e il crollo di tutte le qualita umane \t i personaggi si legano sempre + al discorso indiretto libero (Joyce, Woolf) \\
Le figure sono rese ambigue \t le figure hanno come una maschera, che ricordano "I 6 personaggi in cerca di autore" di Pirandello \\
Personaggi vivono l ipocrisia valoriale, e l assenza di riscatto \t e cercano un evasione, ma riprecipitano nel loro carcere \t per esempio la ragazza ha finalmente un appuntamento, ma prima di uscire di casa viene colta dal vomito (la maggiore espressione del non senso) \\
Lei va oltre le sue aspettative + profonde \t nulla ha senso \t il mondo si presenta attraverso l indifferenza, verso ogni aspettativa e ogni realizzazione \t è un mondo in cui le maschere soffocano le persone \\
La felicita di essere è schiacciata dalla constatazione di non poter essere \t l'aspetto + riv. è il linguaggio, che arriva ai misteri dell esistenza \t le espressioni che noi troviamo, si avvicinano ai quadri di De Chirico (in cui la spersonalizzazione domina) \\
Per lei varcare quella soglia significava togliersi la maschera \t i soliloqui rappresentano una visione dolorosa degli stati d'animo

\ss[Agostino]
ROmanzo fonamentale, un ragazzino che ha una mancata comunicazione con la madre e un senso di ineguatezza con i suoi coetanei \t è vive una omosessualita \\
Simile all isola di arturo, con una figura che avvia al sentimento e crea confusione dei ?? \t qua c'è societa che individua in lui il capro espiatorio, che quell iniziazione è da lui sentita \\
Viene mostrata una seduzione distorta (lui è giovane) \t viene scardinata l ipocrisia della sessualita \t una figura + adulta rappresenta la madre-amante, con una confusione di ruoli che descrive la confusione di una societa \\
Eco pirandelliano si sente \t le maschere sono presenti, l amore è vissuto in modo patetico e umoristico, e il disfacimento del romanzo

\v

Anni 29 e 30 sono fondamentali \t il 30 è l anno della svolta

\ss[Elio Vittorini]
Gli intellettuali sono segnati nell animo dalla guerra \\
Elio vittorini \t è cognato di Quasimodo, e descrive nella prefazione al "Garofano rosso" descrive la situazione dell epoca e il ruolo dell intellettuale, il quale deve medicare e proteggere dal dolore \t questo si vede in:
\bi
  \item "conversazioni in sicilia", 1941 \t delitto maetteotti
  \item "il garofano rosso" \t lettura della parabola dei fatti dell epoca
  \item "erica e i suoi fratelli" \t piaga della prostituzione, mondo post bellico  \t desolazione letta con l animo dell autore \t la ragazza viene mercificata nella prostituzione
  \item "uomini e no"
\ei
Importante \t nella sua immagine di letteratura, in "Uomini e no", c'è constatazoine che ci sono uomini che vivono come uomini, e altri che no \t negazione della dignita dell uomo \\
Vittorini propone un dolore che deve essere schermato, e la letteratura serve a questo (ma non è consolatoria, protegge dal dolore) \\
Parabola del politecnico

\v

I romanzi pongono di fronte anche alla questione della fruizione del romanzo, romanzo come slancio verso una ??? \\
Industrializzazione, ceto proletario \t occasione per leggere, che diventa l affermazione \t la diffusione dei libri \t ritratto che fa Vasco Pratolini, "Cronache di povere amanti" = umanita palpitante, che paga scontri tra fascisti e antifascisti \\
Sofferenza porta pero in questo romanzo alla necessita di leggere, in cui i dialoghi sono presi dal mondo reale \t dalle citta, dalle fabbriche e da queste figure semplici \t che manzoni definirebbe con le "genti meccaniche" \\
Neorealismo si vedono nel romanzo e nel cinema \t Cesare Pavese molto presente
\end{document}
